% =====================================================================================
%  Chapter 9 — Conclusion
%  Thesis: Comparison of Control Strategies for Interaction-Force Aware Multirotor Teams
%
%  DRAFT SKELETON, 27 September 2026.
%  Sources: Section~\ref{sec:intro-rq}, Chapter~\ref{ch:discussion}, project workflow docs.
%
%  HONESTY NOTE. Headline answers per RQ are placeholders until results chapters are filled.
%  The future-work section contains real, already-planned items (not fabricated outcomes).
%
%  Citation baseline: 2026-08-27 (see docs/thesis/CITATION_BASELINE.md).
% =====================================================================================

\chapter{Conclusion}
\label{ch:conclusion}

% -------------------------------------------------------------------------------------
\section{Summary}
\label{sec:conclusion-summary}

This thesis set out to compare interaction-force compensation strategies for tight multirotor
formation flight on a common platform, trajectory set and logging pipeline, holding reactive and
predictive families under the same residual definition and null-flight measurement floor. Strategies~0--2
are implemented, and Strategy~3 is specified but not implemented, in Chapters~\ref{ch:model}--\ref{ch:control}; the experimental
campaign is described in Chapter~\ref{ch:setup} and reported in Chapters~\ref{ch:results2}
and~\ref{ch:results3}. The comparison is deliberately conditional: the useful outcome is guidance
on when to pay for prediction, measurement, or neither, not a claim of universal superiority.

% -------------------------------------------------------------------------------------
\section{Answers to the research questions}
\label{sec:conclusion-rq}

\todo[color=red!25]{PLACEHOLDER (RQ1): One-sentence headline answer --- reactive vs predictive vs
baseline under identical conditions, once Section~\ref{sec:res2-tracking} is populated.}

\todo[color=red!25]{PLACEHOLDER (RQ2): One-sentence headline --- regime dependence or absence thereof,
referencing Section~\ref{sec:res2-summary}.}

\todo[color=red!25]{PLACEHOLDER (RQ3): One-sentence headline --- transfer penalty magnitude and
whether summed architecture remains justified, referencing Section~\ref{sec:res3-transfer}.}

\todo[color=red!25]{PLACEHOLDER (RQ4): One-sentence headline --- cost vs performance trade-off across
strategies (sensing, compute, data).}

% -------------------------------------------------------------------------------------
\section{Engineering takeaway}
\label{sec:conclusion-engineering}

The engineering reader should leave with a decision rule, not only a ranking. RQ4 frames that rule:
given brushless Crazyflies with RPM sensing and offline-trained predictors, which strategy family
delivers acceptable tracking and residual rejection per unit of tuning effort, flight time for data
collection, and onboard resource? The filled discussion (Section~\ref{sec:disc-rq4}) will state that
rule explicitly; until C.4 completes, no rule is asserted here.

% -------------------------------------------------------------------------------------
\section{Future work}
\label{sec:conclusion-future}

Several extensions are already scoped in project documentation and do not depend on the C.4 outcome:

\begin{itemize}
  \item \textbf{C.1 coverage and C.4 repeats.} Complete the planned dz/speed grid under geometric
    collection on cf5; run the systematic comparison matrix with frozen gains and documented $n$
    per cell.
  \item \textbf{Three-robot scaling.} Fly or analyse B1/B2/B3 under the fallback protocol in
    Section~\ref{sec:res3-limitation}; optional sketches for four or more vehicles remain
    a sketch until the core library and the comparison are finished.
  \item \textbf{Retrained NA-INDI / LINDI reference.} Our notes describe
    a retrain path for the Cobo-Briesewitz hybrid; it is out of the primary compared set but remains
    a candidate extension if hardware budget allows.
  \item \textbf{Residual reinforcement-learning strategies.} Explicitly deferred in the project convention:
    reinforcement-learning residual policies are not part of this comparison and would require a
    separate safety and data budget before any flight claim.
  \item \textbf{Hardware validation of onboard prediction.} SIL and desk parity do not substitute for
    \texttt{rnn.en=1} flight logs on the brushless pair; flash-resident vs upload deployment trades
    reflash cost against lab iteration speed.
\end{itemize}

% -------------------------------------------------------------------------------------
\section{Closing remark}
\label{sec:conclusion-close}

Whether compensation families separate sharply or collapse to ``any compensation helps,'' the thesis
is designed so either outcome is informative. The remaining work is to replace the placeholders in
Chapters~\ref{ch:results2}--\ref{ch:discussion} with measured outcomes from the lab campaign.
